\begin{frame}{Considerações Finais}
\begin{itemize}
\item Implementamos a aquisição de 8 bits com temporização pelo Timer1, armazenamento em RAM e transferência serial para o Octave.
\item A taxa máxima aceita pelo firmware é aproximadamente 32,79 kHz. A escolha da frequência do sinal também depende da quantidade de pontos por rampa e do tamanho do bloco.
\item Nos ensaios principais, o erro residual ficou próximo de uma fração de LSB. A nova captura de 10 Hz apresentou RMSE bruto de 0,393 LSB e residual de 0,339 LSB.
\item As medições de bancada mostraram a importância de verificar a tensão aplicada e distinguir a referência nominal da escala efetiva do conjunto.
\end{itemize}
\pblfonte{Relatório técnico; firmware; capturas e medições de bancada.}
\end{frame}

\begin{frame}{Limitações e Continuidade}
\begin{itemize}
\item Medir a referência de forma independente e registrar a incerteza dos instrumentos antes de corrigir definitivamente a escala em volts.
\item Repetir os ensaios com mais blocos e diferentes relações entre frequência do sinal e taxa de amostragem.
\item Verificar todos os níveis em cada rampa, além do histograma acumulado.
\item Comparar a temporização nominal com uma medição externa. Na nova aquisição, estimamos 10,072 Hz e o osciloscópio indicou 10,01 Hz.
\end{itemize}
\begin{block}{Alcance dos resultados}
As medidas caracterizam o conjunto de bancada nas condições ensaiadas. O ajuste da triangular não separa completamente os erros do ADC, do gerador e dos instrumentos.
\end{block}
\end{frame}
